% RecSys Odyssey repository-backed lecture note.
% Add figures with: \coursefigure{figures/name.svg}{Caption}
\section{MMR balances relevance and redundancy}

\begin{abstract}
Each next item should be useful and add something the slate does not already contain.
\end{abstract}

\subsection{Concept}
Maximal Marginal Relevance greedily selects the next item using its relevance minus its maximum similarity to already selected items. The parameter λ controls the trade-off: λ near one behaves like pure relevance; smaller λ penalizes redundancy more strongly. Similarity may use genre, creator, embeddings or calibrated business groups. MMR is transparent and practical, but it cannot fix poor candidates that retrieval never supplied.

\begin{equation*}
MMR(i) = λ\cdotrel(i) - (1-λ)\cdotmax_j∈S sim(i, j)
\end{equation*}

\subsection{Key terms}
\begin{itemize}
\item \textbf{Redundancy.} Repeated information or intent within the same slate.
\item \textbf{λ (lambda).} The control balancing relevance against diversity in MMR.
\item \textbf{Greedy re-ranking.} Building a slate one locally best item at a time.
\end{itemize}
